\begin{landscape}
\begingroup
\fontsize{10}{13}\selectfont
\fancyhead[L]{\small NACHLASS. STÖRUNGEN DER PALLAS}
\renewcommand{\headrulewidth}{0pt}
\par\medskip
\pdfbookmark[2]{Nachlass VI [10]}{pallas-VI-fragment10}
\begin{center}[10.]\end{center}
\par\noindent\hspace*{\parindent}[Zum Behuf der strengen Integration setze man
\[
u=18\text{\mdseries\AstroSymbols ♃}-7\text{\mdseries\AstroSymbols ⚴}+A,\qquad
A=\text{constans},\qquad
\int n\mathrm{d}t=\alpha\iint\sin u\,\mathrm{d}\text{\mdseries\AstroSymbols ⚴}^{2}.
\]
\SourceLeaf{564}
\par\noindent\hspace*{\parindent}Da nun
\[
\text{\mdseries\AstroSymbols ⚴}=\int n\mathrm{d}t+\varepsilon,\quad
\text{also}\quad\frac{\mathrm{d}\mathrm{d}\text{\mdseries\AstroSymbols ⚴}}{\mathrm{d}t^2}=\frac{\mathrm{d}n}{\mathrm{d}t},
\]
so erhält man für \(u\) die Gleichung
\[
\frac{\mathrm{d}\mathrm{d}u}{nn\mathrm{d}t^2}=-7\alpha\sin u,
\]
also
\[
\left(\frac{\mathrm{d}u}{n\mathrm{d}t}\right)^2=C+14\alpha\cos u.
\]
\par\noindent\hspace*{\parindent}Hieraus folgt, dass, wenn \(C\) positiv und dem Betrage nach grösser als \(14\alpha\), \(\dfrac{\mathrm{d}u}{\mathrm{d}t}\) nicht verschwindet und \(u\) beständig zu- oder abnimmt. Ist dagegen \(C\) positiv und kleiner als der Betrag von \(14\alpha\), oder ist \(C\) negativ, so ist \(u\) eine periodische Function und oscillirt um \(0\) oder \(180^0\), je nachdem \(\alpha\) positiv oder negativ ist. In letzterm Falle stehen also die mittlern Bewegungen im rationalen Verhältniss und das Verhältniss \(\dfrac{n'}{n}=\tfrac{7}{18}\) stellt sich immer wieder her.
\par\noindent\hspace*{\parindent}Bezeichnet man den Werth von \(u\) zu irgend einer Anfangszeit mit \(U\) und den entsprechenden Werth von \(\dfrac{\mathrm{d}u}{\mathrm{d}\text{\mdseries\AstroSymbols ⚴}}=18\dfrac{n'}{n}-7\) mit \(\mu\), so wird
\[
C=\mu\mu-14\alpha\cos U
\]
und mit den numerischen Werthen
\begin{center}
\begin{tabular}{@{}lll@{}}
\(\alpha=0\overset{\prime\prime}{,}02955\)&\(n'=299\overset{\prime\prime}{,}12817\)&\(18\text{\mdseries\AstroSymbols ♃}-7\text{\mdseries\AstroSymbols ⚴}=125^0 36'\,49''\)\\
\(A=80^0 56'\,57''\)&\(n=769{,}16512\)&\(U=206\quad33\quad46\)
\end{tabular}
\end{center}
ergibt sich
\[
\begin{gathered}
\log14\alpha=4{,}30221-10\\
\log\mu=6{,}29314-10
\end{gathered}
\qquad\log C=4{,}26300-10.
\]
Es tritt also hier der Fall der Rationalität der mittlern Bewegungen wirklich ein.
\par\noindent\hspace*{\parindent}Weiter wird
\[
\left(\frac{\mathrm{d}u}{n\mathrm{d}t}\right)^2=(C+14\alpha)\left\{1-\frac{28\alpha}{C+14\alpha}\sin\tfrac12u^2\right\},
\]
und durch die Substitution
\[
\BareRoot{\frac{28\alpha}{C+14\alpha}}\mathpunct{\cdot}\sin\tfrac12u=\sin\psi
\]
\[
\frac{\mathrm{d}u}{n\mathrm{d}t}=\BareRoot{(C+14\alpha)}\mathpunct{.}\cos\psi
\]
ergibt sich
\[
\frac{\mathrm{d}\psi}{n\mathrm{d}t}=\BareRoot{7\alpha}\mathpunct{.}\BareRoot{\left(1-\frac{C+14\alpha}{28\alpha}\sin\psi^2\right)}
\]
also
\[
n\mathrm{d}t=\frac{\mathrm{d}\psi}{\BareRoot{\left(7\alpha-\dfrac{C+14\alpha}{4}\sin\psi^2\right)}}\mathpunct{\cdot}
\]
\par\noindent\hspace*{\parindent}Das Maximum (resp. Minimum) von \(\dfrac{\mathrm{d}u}{n\mathrm{d}t}\) ist
\[
=\BareRoot{(C+14\alpha)}=\BareRoot{\left(28\alpha\sin\frac{U^2}{2}+\mu\mu\right)}.
\]
\SourceLeaf{565}
\par\noindent\hspace*{\parindent}Hienach oscillirt die mittlere Bewegung der Pallas zwischen den Werthen:
\[
\frac{18}{7}n'\pm0\overset{\prime\prime}{,}2153.
\]
\par\noindent\hspace*{\parindent}Das Verhältniss der Periode dieser Oscillation zu einem Pallasumlauf ergibt das Integral, von \(0\) bis \(\dfrac{\pi}{2}\) genommen:
\[
\begin{aligned}
\frac{2}{\pi}\int\frac{\mathrm{d}\psi}{\BareRoot{\left(7\alpha-\dfrac{C+14\alpha}{4}\sin\psi^2\right)}}
&=\frac{1}{M\left(\BareRoot{7\alpha},\BareRoot{\dfrac{14\alpha-C}{4}}\right)}\\[.8ex]
&=\frac{1}{M(\BareRoot{7\alpha},\BareRoot{(7\alpha\cos\tfrac12U^2-\tfrac14\mu\mu)})}\mathpunct{\raisebox{.35ex}{,}}
\end{aligned}
\]
d. i. 1894 Pallasumläufe \(=\) 737 Jupitersumläufe.
\par\noindent\hspace*{\parindent}Das Maximum und Minimum von \(u\) selbst ergibt sich zu:
\[
\pm156^0 1'.\text{]}
\]
\clearpage
\endgroup
\end{landscape}
